% from ega2-1-fr.tex
\oldpage[II]{171}
\begin{equation}
\label{II.8.5.3.3-fr}
  \xymatrix{
    G(\varphi)\ar[r]^{\operatorname{Proj}(\varphi)}\ar[d]_{j'} &
    X\ar[d]^{j}\\
    G(\widehat{\varphi})\ar[r]_{\widehat{\Phi}} & \widehat{C}
  }
\tag{8.5.3.3}
\end{equation}
est commutatif.
\end{env}

\begin{proposition}[8.5.4]
\label{II.8.5.4-fr}
Si $\widehat{E}$ (resp. $\widehat{E'}$) est le cône épointé projectif défini
par $\mathcal{S}$ (resp. $\mathcal{S}'$), on a
$\widehat{\Phi}^{-1}(\widehat{E})\subset\widehat{E'}$~; en outre, si
$p:\widehat{E}\to X$ et $p':\widehat{E'}\to X'$ sont les rétractions
canoniques, on a
$p'(\widehat{\Phi}^{-1}(\widehat{E}))\subset G(\varphi)$, et le diagramme
\begin{equation}
\label{II.8.5.4.1-fr}
  \xymatrix{
    \widehat{\Phi}^{-1}(\widehat{E})
      \ar[r]^{\widehat{\Phi}}\ar[d]_{p'} &
    \widehat{E}\ar[d]^{p}\\
    G(\varphi)\ar[r]_{\operatorname{Proj}(\varphi)} & X
  }
\tag{8.5.4.1}
\end{equation}
est commutatif. Si $\operatorname{Proj}(\varphi)$ est partout défini, il en
est de même de $\widehat{\Phi}$, et on a
$\widehat{\Phi}^{-1}(\widehat{E})=\widehat{E'}$.
\end{proposition}

\begin{proof}
La première assertion résulte de la commutativité des diagrammes
(\hyperref[II.8.5.1.3-fr]{8.5.1.3}) et
(\hyperref[II.8.5.3.2-fr]{8.5.3.2}), et les deux suivantes de la définition
des rétractions canoniques (\hyperref[II.8.3.5-fr]{8.3.5}) et de la définition
de $\widehat{\varphi}$. D'autre part, pour voir que $\widehat{\Phi}$ est
partout défini lorsqu'il en est ainsi de $\operatorname{Proj}(\varphi)$, on
peut se borner à considérer le cas où $Y=\operatorname{Spec}(A)$ et
$Y'=\operatorname{Spec}(A')$ sont affines,
$\mathcal{S}=\widetilde{S}$, $\mathcal{S}'=\widetilde{S'}$~; l'hypothèse est
que lorsque $f$ parcourt l'ensemble des éléments homogènes de $S_+$, l'idéal
engendré dans $S'_+$ par les $\varphi(f)$ a pour racine dans $S'_+$ l'idéal
$S'_+$ lui-même~; on en conclut aussitôt que la racine dans $(S'[z])_+$ de
l'idéal engendré par $z$ et les $\varphi(f)$ est $(S'[z])_+$ lui-même, d'où
notre assertion~; cela prouve en même temps que $\widehat{E'}$ est réunion des
$\widehat{C'}_{\varphi(f)}$, donc égal à
$\widehat{\Phi}^{-1}(\widehat{E})$.
\end{proof}

\begin{corollary}[8.5.5]
\label{II.8.5.5-fr}
Lorsque $\operatorname{Proj}(\varphi)$ est partout défini, l'image réciproque
par $\widehat{\Phi}$ de l'espace sous-jacent au lieu à l'infini (resp. au
préschéma des sommets) de $\widehat{C}$ est l'espace sous-jacent au lieu à
l'infini (resp. au préschéma des sommets) de $\widehat{C'}$.
\end{corollary}

\begin{proof}
Cela résulte aussitôt de (\hyperref[II.8.5.4-fr]{8.5.4}) et
(\hyperref[II.8.5.2-fr]{8.5.2}), compte tenu des relations
(\hyperref[II.8.3.4.1-fr]{8.3.4.1}) et
(\hyperref[II.8.3.4.2-fr]{8.3.4.2}).
\end{proof}

\subsection{Un isomorphisme canonique pour les cônes épointés}
\label{subsection:II.8.6-fr}

\begin{env}[8.6.1]
\label{II.8.6.1-fr}
Soient $Y$ un préschéma, $\mathcal{S}$ une $\mathcal{O}_Y$-Algèbre graduée
quasi-cohérente à degrés positifs \emph{telle que
$\mathcal{S}_0=\mathcal{O}_Y$}, $X$ le $Y$-schéma
$\operatorname{Proj}(\mathcal{S})$. Nous allons appliquer les résultats de
(\hyperref[subsection:II.8.5-fr]{8.5}) au cas où on prend $Y'=X$,
$q:X\to Y$ étant le morphisme structural~; soit
\begin{equation}
\label{II.8.6.1.1-fr}
  \mathcal{S}_X=\bigoplus_{n\in\mathbf{Z}}\mathcal{O}_X(n)
\tag{8.6.1.1}
\end{equation}

